\scanpage{164}
% L1–3
donc l'application
\[
\zeta\longmapsto q_\zeta,\quad A*A\longrightarrow
\begin{aligned}&\text{formes quadratiques}\\&\text{sur $E$ qui sont}\\&\text{compatibles à $t$}\end{aligned}
\]
est compatible avec l'application
\[
\begin{aligned}
\deleted{A}L^{\otimes 2}&\longrightarrow L^{\otimes 2}\\
\alpha&\longmapsto-\alpha
\end{aligned}
\]
% L4–7
sur les \deleted{\unsure{modules d'opérateurs|modules d'opération}} \inserted{modules des translations}. Donc il
vaut mieux changer notre convention
et \deleted{\unsure{poser|pour}} prendre $q_\zeta$ tel que
\[
\zeta+q_\zeta(x)=x*x
\]

% L8–10
Ainsi, on a \unsure{un iso.|un isom.} de torseurs
\[
\mathrm{Quadcomp}\underbrace{(E,\epsilon,t)}_{\text{ou }(L,\sigma)}\simeq L*_{\sigma}L
\]

% L11–16
Moyennant quoi \unsure{on} va \unsure{identifier|identifie} les
deux. Ceci posé, pour toute $q\in L*_\sigma L$,
le diviseur quadratique associé n'est
autre que l'image inverse de la section
$q$ de $\check{\mathbb V}(L*_\sigma L)$ par $\mathbb V(L)\mathbin{\text{---}}\mathbb V(L*_\sigma L)$.

% L17–21
\S En termes de coordonnées,
i.e. choisissant une origine dans $\deleted{\unsure{L}}A$,
de sorte que $A$ s'identifie à $L$,
et $t$ est \deleted{\unsure{une}} donné par $\tau\in L$,
\[
\sigma_L(\deleted{\unsure{k}}x)=\tau-x
\]
% archive sheet number: 163
